% conclusions.tex (v14, 2026-10-20): the contributions of this study only; statements already established elsewhere are left to the Discussion.
\section{Conclusions}\label{sec:conclusions}
This study evaluated and mapped active-layer thickness in a site-year design, in which every label, climate input, prediction and score refers to one observation site in one stated year. Three findings are specific to this study. First, a Stefan model recalibrated with local labels and corrected by a gradient-boosting model of its residual reduced the Stefan error by \ndrefLeblkRoabs{} cm in the Lena Delta and \ndrefCablkRoabs{} cm in Canada under spatial-block validation, and the gain was conditional: it required soil and terrain covariates and was largest where the regional Stefan coefficient misrepresented the sites. Second, the physical anchor preserved transferability: without labels of the target region the physics-residual model stayed close to the Stefan model, whereas direct ML exceeded the Stefan error by \ndrefLeregDzabs{} cm in the Lena Delta, and 5--25 local labels were enough to reduce the physics-residual error in the Lena Delta and in external CALM regions. Third, scoring all models and five published products on the same cells and years showed that the accuracy of the ML products follows the location of their training sites; in the Lena Delta they did not represent thin active layers (correlation \nrLeWei{} and \nrLeLiu{} with the observations) and were less accurate than a model trained without any label of the region. The workflow delivers year-labelled 1~km maps for 2011--2024 for three regions, each with an area of applicability and a prediction interval, together with the evaluation protocol and the label-number evidence needed to judge whether, and with how many labels, the method can be applied to another permafrost region.
